\documentclass[10pt,a4paper]{amsart}
\usepackage[margin=19mm]{geometry}
\usepackage{amsmath,amssymb,mathtools}
\usepackage[scr=rsfs,cal=euler]{mathalfa}
\usepackage{xfrac}
\usepackage[unicode,hidelinks,bookmarks=false]{hyperref}
\newcommand{\ie}{i.e.\ }
\newcommand{\cf}{cf.\ }
\newcommand{\resp}{respectively\ }
\newcommand{\Subsection}{\subsection}
\providecommand{\nobreakdash}{\nobreak\hbox{-}}
\setlength{\emergencystretch}{2em}
\multlinegap0pt
\usepackage{amssymb}
\usepackage{tikz}
\usepackage{caption}
\usepackage[all]{xy}
\usepackage[shortlabels]{enumitem}
\usepackage{cancel}
\newtheorem{lemma}{Lemma}
\newcommand{\N}{{\mathbb N}}
\newcommand{\TT}{{\mathbb{T}^2}}
\newcommand{\TTf}{{\mathbb{T}_f^2}}
\newcommand{\cW}{{\mathcal W}}
\newcommand{\eps}{\varepsilon}
\newcommand{\eqdef}{\stackrel{\scriptscriptstyle\textrm def.}{=}}
\newcommand{\tif}{\tilde{f}}
\newcommand{\tix}{\tilde{x}}
\newcommand{\w}{\widetilde{\cW}}
\title{$\mathbf{\textit{U}}$-Gibbs measure rigidity for partially hyperbolic endomorphisms on surfaces}
\author{Marisa Cantarino, Bruno Santiago}
\date{Source: arXiv:2407.03780v2 (CC BY 4.0)}
\begin{document}
\maketitle

\noindent\emph{Source and licence.} $\mathbf{\textit{U}}$-Gibbs measure rigidity for partially hyperbolic endomorphisms on surfaces. Version arXiv:2407.03780v2 (CC BY 4.0), distributed by arXiv under the Creative Commons Attribution 4.0 International licence, http://creativecommons.org/licenses/by/4.0/, as recorded in the arXiv licence metadata for this exact version and shown on the abstract page https://arxiv.org/abs/2407.03780v2. Full licence notice: \texttt{licenses/arxiv-2407.03780-CC-BY-4.0.txt}.\par\smallskip

\noindent\textit{Editorial scope.} Counted unit: the article's lemma lem:quasiisometric (source0.tex lines 1469--1476). The article's complete original proof of it is delivered with it (source0.tex lines 1478--1535, one \texttt{proof} environment of 58 lines). No mathematics is written, repaired or re-derived: the statement and the proof are the article's own lines, in the article's own notation, and no retained line is substituted or re-flowed.  Retained uncounted as the source-local context the proof needs: lines 498--509; lines 945--951; lines 1023--1026; lines 1442--1445.  Nothing was omitted from any retained span, so the package carries no omission record.  No retained line contains a citation, so the wrapper carries no bibliography and no bibliography entry.  No retained line contains a figure environment, an image-inclusion command or drawing code, so no image is delivered and the package declares no asset dependency.  Editorial formatting only, in addition: an AMS \texttt{amsart} preamble that restates the article's own declarations and macros used by the retained lines, namely the environments \texttt{lemma} on separate counters, together with the standard packages those lines need.  The retained environments print in this document as Lemma 1 in order of appearance, whereas the article prints the same items with the numbering of its own full document.  The complete native source of the article as distributed in the arXiv e-print for this exact version is bundled unchanged as \texttt{sources/161816/source0.tex}.
\begin{lemma}
    \label{lem:defideph}
There exists a continuous non-trivial splitting  $T_{\tix}\TTf=E^c(\tix)\oplus E^u(\tix)$ on the inverse limit space, for some choice of background Riemannian metric on $\TT$, such that the continuous functions
\[
\lambda^{*}_{\tix}\eqdef\|D\tilde{f}(\tix)|_{E^{*}}\|,\:\:\:\textrm{for}\:\: * \in \{c,u\},
\]
satisfy the following inequalities:
\begin{enumerate}
\item$\lambda^c_{\tix}<\lambda^u_{\tix}$, for every $\tix\in\TTf$;
\item  $\lambda^u_{\tix}>1$, for  every $\tix\in\TTf$.
\end{enumerate}
\end{lemma}

\begin{lemma}
    \label{lem:crescelyapunovcresce}
There exists a constant $\sigma=\sigma(f)>1$ such that for $\tilde{\mu}$ a.e. $\tix\in\TTf$ and every $v\in E^c(\tix)$ it holds 
\[
e^{n\lambda}\|v\|_{L,\tix}\leq\|D\tif^n(\tix)v\|_{L,\tix_n}\leq e^{n\sigma}\|v\|_{L,\tix},\:\:\:\textrm{for every}\:\:n>0
\]
\end{lemma}

 \begin{equation}
     \label{eq:cocycle}
     \hat{\lambda}^c_{\tix}(m+n) = \hat{\lambda}^c_{\tilde{f}^n(\tilde{x})}(m) \hat{\lambda}^c_{\tix}(n). 
 \end{equation}

\begin{equation}
	\label{eq:tau}
 \tau(n)=\tau(\tix,\tix^u,n,\varepsilon) = \inf \left\{ k\in\N; \; \hat{\lambda}^c_{\tix^u}(k) \frac{\lambda^c_{\tix_{-n}}(n)}{\lambda^u_{\tix_{-n}}(n)} \geq \varepsilon \right\}. 	
\end{equation}

\begin{lemma}
\label{lem:quasiisometric}
There are constants $A >0$, $\Theta > 1$ (depending only on $(f,\mu)$) such that, for $\tilde{\mu}$-almost every $\tix \in \TTf$ and for all $\tix^u \in \w^u(\tix)$, and $m,n \in \N$, we have
\begin{itemize}
	\item $\Theta^{-1}|m - n|-A\leq|\tau(m)-\tau(n)|\leq \Theta|m - n|+A$,
	\item $\Theta^{-1}|m - n|-A\leq|t(m)-t(n)|\leq \Theta|m - n|+A$.
\end{itemize}
\end{lemma}

\begin{proof}
    For a $\tilde{\mu}$-regular point $\tix \in \TTf$, $\tix^u \in \w^u(\tix)$ and $\varepsilon > 0$, consider $\tau(n) = \tau(\tix, \tix^u, n, \varepsilon)$ and $t(n) = t(\tix, \tix^u, n, \varepsilon)$. From \eqref{eq:tau} we have

    $$\hat{\lambda}^c_{\tix^u}(\tau(n)) \frac{\lambda^c_{\tix_{-n}}(n)}{\lambda^u_{\tix_{-n}}(n)} \geq \varepsilon \mbox{, and } \hat{\lambda}^c_{\tix^u}(\tau(n)-1) \frac{\lambda^c_{\tix_{-n}}(n)}{\lambda^u_{\tix_{-n}}(n)} < \varepsilon.$$

    Using equation \eqref{eq:cocycle} and Lemma \ref{lem:crescelyapunovcresce}, the above inequalities imply that
    \begin{equation}
        \label{eq:tau-bounded1}
        \varepsilon \leq \hat{\lambda}^c_{\tix^u}(\tau(n)) \frac{\lambda^c_{\tix_{-n}}(n)}{\lambda^u_{\tix_{-n}}(n)} = \hat{\lambda}^c_{\tilde{f}^{\tau(n) - 1}(\tix^u)}(1) \hat{\lambda}^c_{\tix^u}(\tau(n) - 1) \frac{\lambda^c_{\tix_{-n}}(n)}{\lambda^u_{\tix_{-n}}(n)} < e^\sigma \varepsilon,
    \end{equation}
    for all $n \in \mathbb{N}$.

    By compactness and Lemma~\ref{lem:defideph}, one easily deduces that there are $\chi_1^d < \chi_2^d < 0$ such that $$e^{\chi_1^dn} < d^n_{\tix} \eqdef \frac{\lambda^c_{\tix_{-n}}(n)}{\lambda^u_{\tix_{-n}}(n)} < e^{\chi_2^dn}.$$
    
    By equation \eqref{eq:tau-bounded1} we have, given $m \in \mathbb{N}$, $m \geq n$, that
    $$\varepsilon e^{\chi_1^d(m-n)} \leq \hat{\lambda}^c_{\tix^u}(\tau(n)) \underbrace{d^{m-n}_{\tilde{f}^{n}(\tix)}d^{n}_{\tix}}_{d^{m}_{\tix}} < e^{\chi_2^d(m-n)} e^\sigma \varepsilon.$$
    
    Thus,
    \begin{equation}
        \label{eq:tau-bounded2}
        \varepsilon e^{\chi_1^d(m-n)} e^{\lambda(\tau(m) - \tau(n))} \leq \hat{\lambda}^c_{\tix^u}(\tau(m)) d^{n+m}_{\tix} < e^{\chi_2^d(m-n)} e^{\sigma(\tau(m) - \tau(n) + 1)} \varepsilon
    \end{equation}
    is given by using Lemma \ref{lem:crescelyapunovcresce} and the fact that
    $$\hat{\lambda}^c_{\tilde{f}^{\tau(n)}(\tix^u)}(\tau(m) - \tau(n)) = \dfrac{\hat{\lambda}^c_{\tix^u}(\tau(m))}{\hat{\lambda}^c_{\tix^u}(\tau(n))}.$$

    The left hand inequality on equation \eqref{eq:tau-bounded1} and the right hand side of equation \eqref{eq:tau-bounded2} imply that
    $$\varepsilon \leq e^{\chi_2^d(m-n)} e^{\sigma(\tau(m) - \tau(n) + 1)} \varepsilon,$$
    and, by using the remaining inequalities on said equations, we have also
    $$\varepsilon e^{\chi_1^d(m-n)} e^{\lambda(\tau(m) - \tau(n))} < e^\sigma \varepsilon.$$
    We may assume without loss of generality that $\eps$ is small enough so that $e^{\sigma}\eps<1$. Then, by taking the logarithm on the last two inequalities, we have
    \begin{align*}
        \chi_1^d(m-n) + \lambda(\tau(m) - \tau(n) - 1) &< 0 <  \chi_2^d(m-n) + \sigma(\tau(m) - \tau(n) + 1)\\
        \implies \dfrac{\chi_1^d}{\lambda} (m-n) - 1 &< - \tau(m) + \tau(n) < \dfrac{\chi_2^d}{\sigma} (m-n) + 1,
    \end{align*}
    which gives us the bound for $\tau$.

    A similar argument gives us the bound for $t$. Indeed, the definition of $t(n)$ implies that
    \begin{equation}
        \label{eq:t-bounded}
        1 \leq \dfrac{\hat{\lambda}^c_{\tix}(t(n))}{\hat{\lambda}^c_{\tix^u}(\tau(n))} < e^\sigma.
    \end{equation}
    for all $n \in \mathbb{N}$.
    
    Lemma \ref{lem:crescelyapunovcresce} gives us that
    $$e^{\lambda (t(m) - t(n)) - \sigma (\tau(m) - \tau(n))} < \dfrac{\hat{\lambda}^c_{\tilde{f}^{t(n)}(\tix)}(t(m) - t(n))}{\hat{\lambda}^c_{\tilde{f}^{\tau(n)}(\tix^u)}(\tau(m) - \tau(n))} < e^{\sigma (t(m) - t(n)) - \lambda (\tau(m) - \tau(n))},$$
    which in turn implies
    $$e^{\lambda (t(m) - t(n)) - \sigma (\tau(m) - \tau(n))} < \dfrac{\hat{\lambda}^c_{\tix}(t(m))}{\hat{\lambda}^c_{\tix^u}(\tau(m))} < e^{\sigma (t(m) - t(n) + 1) - \lambda (\tau(m) - \tau(n))},$$
    by equations \eqref{eq:cocycle} and \eqref{eq:t-bounded}.

    The inequalities from both the above equation and \eqref{eq:t-bounded} imply that
    $$e^{\lambda (t(m) - t(n)) - \sigma (\tau(m) - \tau(n) + 1)} < 1 < e^{\sigma (t(m) - t(n) + 1) - \lambda (\tau(m) - \tau(n))},$$
    and then, by taking the logarithm, we have
    \begin{align*}
        \lambda (t(m) - t(n)) - \sigma (\tau(m) - \tau(n) + 1) &< 0 < \sigma (t(m) - t(n) + 1) - \lambda (\tau(m) - \tau(n))\\
        \implies - \dfrac{\sigma}{\lambda} (\tau(m) - \tau(n) + 1) &< - t(m) + t(n) <  1 - \dfrac{\lambda}{\sigma} (\tau(m) - \tau(n)),
    \end{align*}
    which, using the bound for $\tau$, gives us the bound for $t$.
\end{proof}

\end{document}
